% #############################################################################
% This is Chapter 3
% !TEX root = ../main.tex
% #############################################################################
% Change the Name of the Chapter i the following line
\fancychapter{Related Work}
% The following line allows to ref this chapter
\label{chap:related_work}

This chapter reviews previous work that is most directly connected to the
tasks studied in this research. \Cref{sec:rw-portuguese-variety-identification}
covers EP/BP identification, while
\Cref{sec:rw-portuguese-variety-translation} reviews work on rewriting between
the two varieties. \Cref{sec:rw-varieties-beyond-portuguese} broadens the
discussion to language processing beyond Portuguese, including work on
dialectal and regional settings. \Cref{sec:rw-sequential-adaptation-mt}
reviews staged adaptation for machine translation, and
\Cref{sec:rw-reward-based-ft-translation} covers reward-based adaptation for
translation. \Cref{sec:rw-overview} identifies the research gap.

\section{European and Brazilian Portuguese Identification}
\label{sec:rw-portuguese-variety-identification}

Early work on EP/BP identification showed that simple surface features are
already highly informative for this distinction.
\citet{DBLP:conf/konvens/ZampieriG12} proposed one of the first automatic
systems for discriminating European and Brazilian Portuguese, using character
\(n\)-grams and word \(n\)-gram language models. Their results indicate that orthographic and lexical cues can be highly
informative for EP/BP identification. More broadly, work on document language
identification cautions against treating closely related varieties as an easy
identification problem \citep{4052828}.

\citet{7839597} tested the same distinction on Twitter, where spelling
variation, limited context, and informal writing make EP/BP identification less
controlled than in edited corpora. Their use of smoothed \(n\)-gram models and
an ensemble strategy shows that lexical and character-level evidence remains
useful in social-media settings, although its reliability depends on the amount
and quality of the available context. \citet{preda-etal-2024-across} introduce
stronger experimental controls through a Portuguese dialect-identification
dataset derived from TED2020, a parallel corpus containing aligned European and
Brazilian Portuguese versions of the same talks. The authors compare handcrafted
linguistic features, word \(n\)-gram classifiers with optional part-of-speech
(POS) information, adaptive Naive Bayes variants, and LoRA-adapted models. Their
evaluation varies input length, from sentence-level segments to full
transcripts, and tests cross-dataset robustness on out-of-domain news sources,
which makes text granularity and domain shift part of the evaluation rather
than minor experimental settings. The LoRA-adapted Albertina classifiers obtain
the strongest in-domain results on TED2020, while a logistic-regression bigram
model performs best in the joint cross-dataset evaluation.

The reliability of these comparisons depends heavily on dataset construction,
particularly when labels are inferred from metadata such as source country,
domain, or Uniform Resource Locator (URL) \citep{zampieri-etal-2014-report}. For very similar varieties,
metadata can be weak supervision, as a text collected from a given country may
contain few explicit markers of that variety, or may have been written in a
different variety altogether. \citet{zampieri-etal-2024-language} address this
issue with Discriminating Between Similar Languages -- True Labels (DSL-TL), a human-annotated multilingual dataset for language-variety
identification that includes European and Brazilian Portuguese, highlighting that high scores on silver-labeled corpora do not necessarily
imply robust variety discrimination.

More recently, \citet{sousa2025enhancingportuguesevarietyidentification}
compiled PtBrVarId, a cross-domain dataset covering journalistic, legal,
political, web, social-media, and literary texts, and evaluated both
\(n\)-gram-based baselines and Transformer-based classifiers. Their best
systems fine-tune BERTimbau \citep{souza2020bertimbau} and use a two-step
cross-domain protocol with partial delexicalization, masking named entities
and selected part-of-speech categories to reduce overfitting to domain-specific
lexical cues. The delexicalized BERT model is their strongest system on both DSL-TL and FRMT.
Their FRMT identification results provide one of the closest published
comparisons for this dissertation's identification task.

\section{Rewriting Between European and Brazilian Portuguese}
\label{sec:rw-portuguese-variety-translation}

\citet{marujo-etal-2011-bp2ep} studied the specific task of converting BP texts
into EP, rather than generating Portuguese text more generally. Their
rule-based and statistical setting was used to adapt BP resources for
EP-oriented machine translation, and the adapted resources improved translation
quality compared with using the same material without the BP-to-EP conversion
step. A complementary linguistic perspective is provided by
\citet{barreiro-mota-2018-paraphrastic}, who extracted paraphrastic
correspondences between both Portuguese variants from a parallel corpus. Their analysis shows that some EP/BP correspondences involve multiword
expressions and syntactic changes, so rewriting is not always word replacement.

In the neural setting, \citet{costa-jussa-etal-2018-neural} presented a system
for translating between EP and BP in both directions and compared it with
phrase-based statistical machine translation. Their recurrent encoder-decoder model improved automatic scores and was also
preferred in human evaluation, indicating that neural models can learn
cross-variety transformations that challenge phrase-based methods.

\citet{sanches2024brazilianportugueseeuropeanportuguese} bring this task closer
to current pretrained systems by focusing on BP-to-EP translation with modern
encoder-decoder and decoder-only architectures. The authors provide a
comparison point for the rewriting component of this work by fine-tuning
multilingual encoder-decoder models, such as M2M100 and mBART50, on parallel
data from subtitles and TED talks. They introduce the manually curated Golden
Collection benchmark and compare the fine-tuned models with ChatGPT 3.5 using
both automatic metrics and human evaluation. Their study finds that the best fine-tuned system
generalizes well on test data close to its training domains, while ChatGPT 3.5
is stronger on the manually curated Golden Collection. Copying the input is also a
very strong baseline against the manual references in that collection,
reflecting the conservative nature of many valid EP/BP conversions.

\section{Language Processing Beyond Portuguese}
\label{sec:rw-varieties-beyond-portuguese}

Language- and variety-identification systems can affect not only evaluation but
also the data available for later model training \citep{zampieri2020natural}.
OpenLID \citep{burchell-etal-2023-openlid} and OpenLID-v3
\citep{fedorova-etal-2026-openlidv3} provide multilingual
language-identification resources for data curation, where predictions can
determine whether candidate web texts are kept, discarded, or routed to a
language-specific subset. FastSpell \citep{banon-etal-2024-fastspell}
specifically addresses confusable languages by combining an initial
language-identifier prediction with dictionary-based verification.
\citet{caswell-etal-2020-language-id-wild} show why this distinction matters in
practice: strong benchmark F1 does not necessarily produce clean web-filtered
corpora under domain mismatch, class imbalance, and language similarity.
Identification errors can therefore propagate into the composition of the
corpora used for subsequent adaptation.

Related work on generation shows that national and regional varieties can also
be represented explicitly as output targets rather than collapsed under a
single language label. \citet{lakew-etal-2018-neural} train multilingual NMT
systems from English into specific national varieties and closely related
standardized languages, including European and Brazilian Portuguese, and find
that treating the varieties as distinct targets improves translation
consistency in their setting. The task remains translation from English rather
than rewriting directly between two Portuguese varieties.
\citet{sousa2025tradutorbuildingvarietyspecific} provide a more recent
Portuguese example with PTradutor, an open-source English-to-EP translation
model built with retro-translated data. Its improvement over generic
open-source Portuguese translation systems further shows that targeting the
desired national variety can matter during data construction and model
adaptation.

Broader dialect and language-variety benchmarks show that the usefulness of
different model families depends on the task being evaluated.
\citet{faisal-etal-2024-dialectbench} introduce DIALECTBENCH, a large-scale
benchmark covering 281 varieties and 10 text-level tasks. Building on this
benchmark, \citet{faisal-anastasopoulos-2025-testing} compare instruction-tuned
LLMs with supervised encoder-based models and find that LLMs can be competitive
on some comprehension-style tasks, while multi-class identification often still
benefits from task-specific supervised fine-tuning. The two studies therefore
provide a benchmark and a subsequent model comparison within the same broader
evaluation setting, rather than independent evidence about unrelated
dialect-processing problems.

Work that combines identification with generation is particularly relevant to
the shared formulation studied in this dissertation.
\citet{khered-etal-2025-multi} combine Arabic dialect identification and
dialect-to-Modern-Standard-Arabic translation within a multi-task
encoder-decoder model. Their experiments report improved translation over the
evaluated baselines, but the study changes both the shared multi-task setup and
the training data, so the gain supports the combined experimental design
without isolating multi-task learning itself as the cause.
\citet{sharma2026indicdialect} introduce INDIC-DIALECT, a human-curated corpus
and multi-task benchmark for Hindi and Odia dialects that includes dialect
classification, multiple-choice question answering, and machine translation.
Their experiments find weak zero-shot LLM classification, substantially
stronger fine-tuned Indic models, and translation results in which hybrid
rule-based and neural systems perform best, with the preferred strategy
depending on conversion direction. Together with the preceding work, these
studies place the present EP/BP setting within a broader literature in which
variety information affects corpus selection, generation targets, and joint
task modeling.

\section{Staged Adaptation for Machine Translation}
\label{sec:rw-sequential-adaptation-mt}

Classical NMT domain adaptation trains models on general-domain bitext before
adapting them to in-domain data, assigning broad and specialized supervision
different roles across successive phases.
\citet{xu2024paradigmshiftmachinetranslation} introduce the Advanced Language
Model-based trAnslator (ALMA), a decoder-only MT system adapted first with
monolingual data and then with a smaller set of high-quality parallel data.
Their ablations show benefits from the monolingual adaptation stage and from
high-quality parallel supervision, without isolating the effect of stage
order.

\citet{alves2024toweropenmultilingual} use a related staged recipe in Tower,
combining continued pretraining on multilingual monolingual and parallel data
with later fine-tuning on translation-related instructions.
\citet{rei2025towerbridginggenerality} extend this recipe by adding preference
optimization and reinforcement learning under verifiable reward signals, so
later phases introduce objectives that are absent from the earlier supervised
stages. Omnilingual MT
\citep{omnilingualmtteam2026omnilingualmtmachinetranslation} assigns distinct
roles to several training phases at larger scale. Its decoder-only system
starts from LLaMA instruction models, extends the tokenizer for multilingual
coverage, continues pretraining on monolingual language modeling and parallel
translation data, and then uses supervised fine-tuning to recover
instruction-following behavior while specializing the model for MT. Its
encoder-decoder OMT-NLLB system uses decoder training on parallel MT and
monolingual autoencoding, cross-attention warm-up, and final end-to-end
fine-tuning on parallel data. The reported gains from autoencoding, warm-up,
and the final stage show that successive phases can serve different functions
rather than simply extending the total amount of training.

Sequential specialization can also weaken behavior acquired in earlier
training phases. Continual-learning studies examine this retention problem
explicitly in NMT
\citep{cao-etal-2021-continual,gu-feng-2020-investigating,gu-etal-2022-continual}.
\citet{shao-feng-2022-overcoming} show that related degradation can also arise
from imbalanced exposure to training examples outside a strict
continual-learning setting. These results make retention across later
adaptation phases an important consideration without treating staged
adaptation and continual learning as the same setting.

\section{Reward-Based Adaptation for Translation}
\label{sec:rw-reward-based-ft-translation}

\citet{nguyen-etal-2017-reinforcement} show that an existing NMT system can be
updated from scalar feedback without receiving a corrected translation for
every generated sentence. Their bandit NMT experiments use simulated human
ratings and optimize corpus-level MT objectives with policy-gradient updates.
\citet{thu2023rlnmt} apply reward-based continuation after supervised training
for low-resource Burmese dialect translation and report improvements over
non-RL baselines. In both studies, reward optimization modifies an already
trained translation system rather than defining the translation task from
scratch.

\citet{he-etal-2024-improving} use Quality Estimation (QE), which predicts
translation quality without requiring a reference translation, as a reward for
MT feedback training. They report translation-quality gains but also observe
overoptimization, where the learned reward continues to improve while actual
translation quality deteriorates. \citet{xu2024advancingtranslationpreference}
train a reward model from human translation preferences and report
improvements over their supervised starting model after Reinforcement Learning
from Human Feedback (RLHF)-based translation optimization, together with
evidence that part of the learned preference signal transfers to translation
directions not used for RLHF training.

\citet{ramos2024finegrainedrewardoptimization} use xCOMET error spans and
severity predictions to replace a single sentence-level score with token-level
rewards. Their experiments report better translation quality than
sentence-level reward baselines and more stable reward progression during
training. \citet{feng2025mtr1zero} combine format constraints with
metric-based translation quality in a mixed reward and show that lexical,
semantic, and mixed metric rewards lead to different translation behavior
during continuation. \citet{pmlr-v235-wang24ay} show that reward scaling,
transformations, and aggregation affect behavior under joint optimization.

Omnilingual MT \citep{omnilingualmtteam2026omnilingualmtmachinetranslation}
reports an optimization failure in which straightforward Group Relative Policy
Optimization (GRPO) continuation after SFT produces low entropy and vanishing
gradients. The revised setup enlarges candidate groups, reintroduces
Kullback--Leibler (KL) regularization, and uses a balanced multi-criterion
reward. \mbox{\citet{liu-etal-2026-mending}} show a complementary failure mode
in multilingual MT, where source-based QE rewards can amplify weaknesses in
the evaluator. These results show that reward-based MT can fail through either
reward design or optimization.

\section{Overview and Research Positioning}
\label{sec:rw-overview}

The closest published experimental comparisons are
\citet{sousa2025enhancingportuguesevarietyidentification} for EP/BP
identification on FRMT and
\citet{sanches2024brazilianportugueseeuropeanportuguese} for BP-to-EP
rewriting on the Golden Collection. Neither study evaluates identification and
bidirectional rewriting across both resources. \citet{sousa2025enhancingportuguesevarietyidentification} provide the most
direct published identification comparison used in \Cref{chap:evaluation},
while \citet{sanches2024brazilianportugueseeuropeanportuguese} provide the most direct published rewriting comparison
and introduce the Golden Collection benchmark.

In Arabic, \citet{khered-etal-2025-multi} combine dialect identification with
dialect-to-Modern-Standard-Arabic translation in a shared encoder-decoder
model. Their translation gains are reported for a setup that also changes the
available training data, so they do not isolate the contribution of multi-task
learning alone. \citet{khered-etal-2025-multi} study one-way dialect-to-standard translation,
whereas this work combines variety identification with bidirectional EP/BP
rewriting and alternates between short label outputs and full-sentence
generation in the same model.

\citet{caswell-etal-2020-language-id-wild} show that language-identification
scores can overstate corpus quality when labels are noisy, classes are
imbalanced, or closely related varieties are mixed. Staged MT studies show
that successive adaptation phases can serve different roles and that later
training can alter earlier behavior
\citep{alves2024toweropenmultilingual,rei2025towerbridginggenerality,omnilingualmtteam2026omnilingualmtmachinetranslation}.
Reward-based MT studies show that evaluator choice, reward composition, and
optimization stability can materially affect the behavior learned during
continuation
\citep{he-etal-2024-improving,feng2025mtr1zero,pmlr-v235-wang24ay}.

Within the literature reviewed in this chapter, no single study combines EP/BP
identification and bidirectional rewriting in a shared model while evaluating
both tasks across FRMT and the Golden Collection. This work studies that
combined setting with encoder-decoder models and examines how encoder-side and
decoder-side control, supervision composition, staged adaptation, and
reward-based continuation affect identification and conservative rewriting.
The rewriting objective is to preserve meaning while making only the changes
needed for the requested Portuguese variety.
